\section{Oriented derivation of the conditional kernel}
\label{app:conditional-kernel}

This appendix supplies the orientation details used in
Proposition~\ref{prop:conditional-law}.  Fix the selected link $e$ and write
$U=U_e$.  For each incident plaquette $p$, choose an initial vertex and
write its holonomy as
\begin{equation}\label{eq:oriented-plaquette}
  U_p=A_pU^{\epsilon_p}B_p,
  \qquad \epsilon_p\in\{-1,+1\},
\end{equation}
where $A_p$ and $B_p$ are ordered products of the other three link variables,
with inverses dictated by orientation.

The heat kernel is a class function.  If $\epsilon_p=+1$, then
\[
  A_pUB_p \quad\hbox{is conjugate to}\quad U B_pA_p,
\]
so the factor equals $p_t(US_p)$ with $S_p=B_pA_p$.  If
$\epsilon_p=-1$, inversion symmetry and cyclic conjugacy give
\[
 p_t(A_pU^{-1}B_p)
 =p_t(B_p^{-1}UA_p^{-1})
 =p_t(UA_p^{-1}B_p^{-1}),
\]
so $S_p=A_p^{-1}B_p^{-1}$.  This proves the staple representation
\eqref{eq:staple-form} for either orientation.  Each $S_p$ is a continuous
function of the conditioned links.

In the full density \eqref{eq:finite-measure}, every nonincident plaquette
factor is independent of $U$ and cancels from the conditional ratio.  Hence
the continuous conditional kernel is
\begin{equation}\label{eq:general-conditional}
 \gamma_{e,t}^{\omega}(\dd U)
 =Z_{e,t}(\omega)^{-1}
 \prod_{p\ni e}p_t\bigl(U S_p(\omega)\bigr)\,\Haar.
\end{equation}
Because all heat kernels are positive, $Z_{e,t}(\omega)>0$.  Joint
continuity on the compact boundary space also makes
$Z_{e,t}(\omega)$ continuous, so \eqref{eq:general-conditional} defines a
regular conditional probability at every boundary assignment, including
the center-valued one.

For the pattern \eqref{eq:one-frustrated-pattern}, substitution into
\eqref{eq:general-conditional} gives
\[
 \prod_{p\ni e}p_t(US_p)=p_t(U)^{m-1}p_t(Uz^{-1}).
\]
Since $z$ is central, $Uz^{-1}=z^{-1}U$, yielding
\eqref{eq:conditional-law}.  The same centrality is what permits a private
opposite link to set a staple without introducing an ordering ambiguity.

Under a vertex gauge transformation, the target plaquette holonomy is
conjugated.  Thus its event
$\dist(\one,U S_q)<\rho$ is gauge invariant even though the conditioning on
individual link coordinates is a coordinate-level operation.
